\documentclass{article}
\usepackage{hyperref}
\usepackage{xcolor}
\usepackage{amsfonts, amsmath, amssymb, amsthm}

\begin{document}

\begin{center}
    {\bfseries\large Journal de bord -- Predoc HEG Genève} \\[0.2cm]
    {\small Lubin Bénéteau}\\[0.2cm]
    {\small 2026 -- 2027}
\end{center}

\section*{Octobre}

\paragraph{Jeudi 1\textsuperscript{er} octobre.} Arrivée à la HEG, récupération du matériel informatique et carte d'accès. Réunion en début d'après midi avec Prof. Weber pour établir les premières consignes. 
\begin{itemize}
    \item Préparation d'un workflow (exploration des données météo $\longrightarrow$ sélection des variables pertinentes (relatives aux vents d'abord, puis celles du soleil et de la température, coordonnées des stations pour création de cartes des vents) $\longrightarrow$ création d'un large dataset $\longrightarrow$ analyse descriptives du dataset.
    \item extraction depuis \href{https://www.meteosuisse.admin.ch/services-et-publications/service/open-data.html#ogd-kw=Vent&ogd-q=&ogd-cat=&ogd-fmt=}{\textcolor{blue}{ici}} et \href{https://github.com/MeteoSwiss/opendata/tree/main/metadata-files}{\textcolor{blue}{là}} des meta paramètres pour une explication et sélection des variables pertinentes pour la production d'énergie éolienne. Code ne Python. 
    \item Création d'un repo Github (\href{https://github.com/LeLubb/Predoc-HEG/tree/main}{\textcolor{blue}{ici}}) pour monitorer mon avancement dans mon travail.
\end{itemize}

\paragraph{Vendredi 2 octobre.} Avancement sur le code
\begin{itemize}
    \item Sélection des variables pertinentes pour le vent (implique lecture sur les types de vents dans la production électrique éolienne).
    \begin{itemize}
        \item Compréhension des mécanismes de création du vent.
        \item Application au cas de la Suisse.
        \item Sélection finale des variables.
    \end{itemize}
    \item Construction d'une base de donnée csv avec python. Utilisation du site Open Data de MétéoSuisse. Rassemblement des différentes stations. Merge de différents types de timming. Merge des stations.
    \item \textcolor{blue}{Il faudra faire de même pour le soleil et la température et ensuite établir une corrélation (ou pas) entre les deux gradeurs}
\end{itemize}



\end{document}
